\newtheorem{definition}{D\'efinition}[chapter]
\newtheorem{theorem}{Th\'eor\`eme}[chapter]
\newtheorem{proposition}{Proposition}[chapter]
\newtheorem{property}{Propri\'et\'e}[chapter]
\newtheorem{corollary}{Corollaire}[chapter]
\newtheorem{problem}{Probl\'eme}[chapter]
\newtheorem{lemma}{Lemme}[chapter]
\newtheorem{claim}{Claim}[chapter]
\newenvironment{proof}{\textit{\textbf{D\'emonstration.}}}{\hfill$\square$\\}

%%% colors
\newcommand{\cDarkRed}[1]{\color[HTML]{660A10} #1}
\newcommand{\cDarkBlue}[1]{\color[HTML]{00197F} #1}
\newcommand{\cDarkGreen}[1]{\color[HTML]{1D660A} #1}
\newcommand{\cGreen}[1]{\color[HTML]{91CF50} #1}
\newcommand{\cGreyBlue}[1]{\color[HTML]{99B4D3} #1}
\newcommand{\cPurple}[1]{\color[HTML]{70309F} #1}
\newcommand{\cLightOrange}[1]{\color[HTML]{FFC000} #1}


%%% references
\newcommand{\refPrePost}[3]{\hyperref[#2]{#1~\vref*{#2}\xspace#3}}
\newcommand{\refPrePostN}[3]{\hyperref[#2]{#1\ref*{#2}\xspace#3}} %%temp
\newcommand{\RefEquation}[1]{\refPrePost{\'equation~}{#1}{}}
\newcommand{\RefEquations}[2]{\refPrePostN{\'equations~}{#1}{}~et~\refPrePost{}{#2}{}}
\newcommand{\RefChapitre}[1]{\refPrePost{chapitre~}{#1}{}}
\newcommand{\RefSection}[1]{\refPrePost{paragraphe~}{#1}{}}
\newcommand{\RefFig}[1]{\refPrePost{Fig.~}{#1}{}}
\newcommand{\RefFigure}[1]{\refPrePost{figure~}{#1}{}}
\newcommand{\RefFigures}[2]{\refPrePost{figures~}{#1}{}~et~\refPrePost{}{#2}{}}
\newcommand{\RefTable}[1]{\refPrePost{table~}{#1}{}}
\newcommand{\RefDefinition}[1]{\refPrePost{d\'efinition~}{#1}{}}
\newcommand{\RefProposition}[1]{\refPrePost{proposition~}{#1}{}}
\newcommand{\RefProperty}[1]{\refPrePost{propri\'et\'e~}{#1}{}}
\newcommand{\RefLemma}[1]{\refPrePost{lemme~}{#1}{}}
\newcommand{\RefTheorem}[1]{\refPrePost{th\'eor\`eme~}{#1}{}}
\newcommand{\RefTheoremFake}[2]{th\'eor\`eme~#1 de \cite{#2}}
\newcommand{\Equ}[1]{\refPrePost{\'equation~}{#1}{}}
\newcommand{\Implies}{\ensuremath{\Rightarrow}}
% such that
\newcommand{\st}{\ensuremath{~|~}}

%%% boring french words
\newcommand{\cad}{c'est-\`a-dire\xspace}
\newcommand{\resp}{resp.\xspace}
\newcommand{\respp}{respectivement\xspace}
\newcommand{\ssi}{si et seulement si\xspace}
\newcommand{\etal}{\emph{et al.}\xspace}
\newcommand{\nauthor}[1]{\textsc{#1}\xspace}
\newcommand{\cauthor}[2]{\textsc{#1}~\cite{#2}\xspace}
\newcommand{\nauthors}[1]{\textsc{#1} \etal\xspace}
\newcommand{\cauthors}[2]{\textsc{#1} \etal~\cite{#2}\xspace}
\newcommand{\ComputerGraphics}{\emph{informatique graphique\xspace}}
\newcommand{\GeometryProcessing}{\emph{traitement g\'eom\'etrique\xspace}}

%%% spaces
\newcommand{\Z}{\ensuremath{\mathbb{Z}}}
\newcommand{\R}{\ensuremath{\mathbb{R}}}
\newcommand{\Vect}[1]{\ensuremath{\overrightarrow{#1}}}
\newcommand{\vz}{\ensuremath{\mathbf{z}}}
\newcommand{\vy}{\ensuremath{\mathbf{y}}}
\newcommand{\vx}{\ensuremath{\mathbf{x}}}
\newcommand{\vt}{\ensuremath{\mathbf{t}}}
\newcommand{\vxH}{\ensuremath{\hat{\mathbf{x}}}}
\newcommand{\vp}{\ensuremath{\mathbf{p}}}
\newcommand{\vu}{\ensuremath{\mathbf{u}}}
\newcommand{\vv}{\ensuremath{\mathbf{v}}}
\newcommand{\vw}{\ensuremath{\mathbf{w}}}
\newcommand{\vnu}{\ensuremath{\mathbf{\nu}}}
\newcommand{\vn}{\ensuremath{\mathbf{n}}}
\newcommand{\vZ}{\ensuremath{\mathbf{Z}}}
\newcommand{\vN}{\ensuremath{\mathbf{N}}}
\newcommand{\NU}{\ensuremath{\mathcal{U}}}

%% Shape notations
\newcommand{\Shape}{\ensuremath{X}}
\newcommand{\DigShape}{\ensuremath{Z}}
\newcommand{\Boundary}[1]{\ensuremath{\partial #1}}
\newcommand{\DigBoundary}[1]{\ensuremath{Bd(#1)}}
\newcommand{\Shapes}{\ensuremath{\mathbb{X}}}

% Any digitization
\newcommand{\AnyDig}{\ensuremath{\mathtt{D}}}
\newcommand{\AnyDigF}[2]{\ensuremath{\AnyDig_{#2}(#1)}}
\newcommand{\AnyDSh}{\AnyDigF{\Shape}{h}}

% Gauss digitization
\newcommand{\Dig}{\ensuremath{\mathtt{G}}}
\newcommand{\DigF}[2]{\ensuremath{\Dig_{#2}(#1)}}
\newcommand{\DSh}{\DigF{\Shape}{h}}

% Inner and Outer Jordan digitization
\newcommand{\JI}{\ensuremath{\mathtt{J}^-}}
\newcommand{\JO}{\ensuremath{\mathtt{J}^+}}
\newcommand{\JIF}[2]{\ensuremath{\JI_{#2}(#1)}}
\newcommand{\JOF}[2]{\ensuremath{\JO_{#2}(#1)}}
% Body of a digitization
\newcommand{\Body}[2]{\ensuremath{\lbrack #1 \rbrack_{#2}}}

% p1,...pn moment
\newcommand{\Mom}[1]{\ensuremath{m^{#1}}}
\newcommand{\DMom}[2]{\ensuremath{\hat{m}_{#2}^{#1}}}
% covariance matrix
\newcommand{\Cov}[0]{\ensuremath{\mathcal{V}}}
\newcommand{\DCov}[1]{\ensuremath{\Cov_{#1}}}


%\newcommand{\Dig}{\ensuremath{D}}
\newcommand{\BT}[1]{\ensuremath{\partial #1}}
\newcommand{\Bd}[1]{\ensuremath{\partial #1}}
\newcommand{\dS}{\BT X}
\newcommand{\dD}{\BT \DSh}

\newcommand{\Proj}{\pi_h^\Shape}
\newcommand{\Card}{\mathrm{Card}}

\newcommand{\q}{\ensuremath{\kappa}}
%\newcommand{\estq}{\ensuremath{\hat{\kappa}}}
\newcommand{\estq}[1]{\ensuremath{\hat{\kappa}_{#1}}}
\newcommand{\speed}[1]{\ensuremath{\tau_{#1}}}
\newcommand{\eabs}{\ensuremath{\epsilon_{abs}}}
\newcommand{\erel}{\ensuremath{\epsilon_{rel}}}
\newcommand{\eavgabs}{\ensuremath{\overline{\epsilon}_{abs}}}
\newcommand{\eavgrel}{\ensuremath{\overline{\epsilon}_{rel}}}

%%Edit macros
\newcommand{\com}[1]{\textcolor{red}{\emph{#1}}}
\newcommand{\comDav}[1]{\textcolor{blue}{ #1}}
\newcommand{\comJeremy}[1]{\textcolor[rgb]{0.26,0.73,0.06}{\emph{#1}}}
\newcommand{\todoJeremy}[1]{\todo[color=green!40,author=Jeremy]{#1}}
\newcommand{\todoInlineJeremy}[1]{\todo[inline,color=green!40,author=Jeremy]{#1}}
\newcommand{\comJaco}[1]{\textcolor{magenta}{#1}}

\newcommand{\MCard}{\ensuremath{\mathrm{Card}}}
\newcommand{\Area}{\ensuremath{\mathrm{Area}}}
\newcommand{\Vol}{\ensuremath{\mathrm{Vol}}}
\newcommand{\AreaC}[0]{\ensuremath{\widehat{\Area}}}
\newcommand{\VolC}[0]{\ensuremath{\widehat{\Vol}}}
% \newcommand{\EqDef}{\ensuremath{\stackrel{def}{=}}}
\newcommand{\EqDef}{\ensuremath{:=}}

%% Area/Volume notation
\newcommand{\AR}{\ensuremath{A(R,x)}}

%% Curvature notations
\newcommand{\Curv}{\ensuremath{\kappa}}
\newcommand{\MeanCurv}{\ensuremath{H}}
\newcommand{\GaussCurv}{\ensuremath{K}}
\newcommand{\PrincCurv}[1]{\ensuremath{\kappa_{#1}}}
\newcommand{\PrincDir}[1]{\ensuremath{\vw_{#1}}}
\newcommand{\NormalDir}{\ensuremath{\vn}}

%% Pottmann curvature estimators
\newcommand{\CurvT}[1]{\ensuremath{\tilde{\Curv}^{#1}}}
\newcommand{\MeanCurvT}[1]{\ensuremath{\tilde{\MeanCurv}^{#1}}}
\newcommand{\GaussCurvT}[1]{\ensuremath{\tilde{\GaussCurv}^{#1}}}
\newcommand{\PrincCurvT}[2]{\ensuremath{\tilde{\PrincCurv{}}_{#1}^{#2}}}
\newcommand{\PrincDirT}[2]{\ensuremath{\tilde{\PrincDir{}}_{#1}^{#2}}}
\newcommand{\NormalDirT}[1]{\ensuremath{\tilde{\NormalDir{}}^{#1}}}

%% II curvature estimators
\newcommand{\CurvH}[1]{\ensuremath{\hat{\Curv}^{#1}}}
\newcommand{\MeanCurvH}[1]{\ensuremath{\hat{\MeanCurv}^{#1}}}
\newcommand{\GaussCurvH}[1]{\ensuremath{\hat{\GaussCurv}^{#1}}}
\newcommand{\PrincCurvH}[2]{\ensuremath{\hat{\PrincCurv{}}_{#1}^{#2}}}
\newcommand{\PrincDirH}[2]{\ensuremath{\hat{\PrincDir{}}_{#1}^{#2}}}
\newcommand{\NormalDirH}[1]{\ensuremath{\hat{\NormalDir{}}^{#1}}}

\newcommand{\BRShape}{\ensuremath{B_{R}(x) \cap \Shape}}

%% II param-free curvature estimators
\newcommand{\CurvHS}[1]{\ensuremath{\CurvH{#1}^{*}}}
\newcommand{\MeanCurvHS}[1]{\ensuremath{\MeanCurvH{#1}^{*}}}

%% II feature estimators
\newcommand{\FeatDD}{\ensuremath{G_{\Shape,x}}}
\newcommand{\FeatTD}{\ensuremath{\mathcal{G}_{\Shape,x}}}
\newcommand{\DigFeatDD}{\ensuremath{G_{\DigShape,p}}}
\newcommand{\DigFeatTD}{\ensuremath{\mathcal{G}_{\DigShape,p}}}
\newcommand{\CurvMax}{\ensuremath{\Curv_{max}}}
\newcommand{\MeanCurvMax}{\ensuremath{\MeanCurv_{max}}}

\newcommand{\anglais}[1]{\emph{#1}\xspace}
\newcommand{\featflat}{\textsc{Flat}\xspace}
\newcommand{\featsmooth}{\textsc{Smooth}\xspace}
\newcommand{\featedge}{\textsc{Edge}\xspace}
\newcommand{\Feature}{\textsc{Feature}\xspace}
\newcommand{\NonFeature}{\textsc{Non-feature}\xspace}
\newcommand{\VCMM}{\anglais{Voronoi Covariance Measure}\xspace}
\newcommand{\VCM}{\anglais{VCM}\xspace}


% \newcommand{\CurvT}[1]{\ensuremath{\tilde{\kappa}_{#1}}}
% \newcommand{\MeanCurvT}[1]{\ensuremath{\tilde{H}_{#1}}}
% \newcommand{\PCurvHat}[1]{\ensuremath{\hat{\Curv}_{#1}}}
% \newcommand{\DCurvHat}[1]{\ensuremath{\hat{\Curv}^D_{#1}}}
% \newcommand{\PMeanCurvHat}[1]{\ensuremath{\hat{H}_{#1}}}

% Digital frontier at step $h$.
\newcommand{\DFr}[1]{\ensuremath{\Delta_{#1}}}
\newcommand{\hBd}[1]{\ensuremath{\partial_{#1}}}

%Ball of radius R (#1) centered on x (#2).
\newcommand{\Ball}[2]{\ensuremath{B_{#1}(#2)}}
%Positive half-Ball of radius R (#1) centered on x (#2).
\newcommand{\HalfBall}[2]{\ensuremath{B^+_{#1}(#2)}}
% Vector #1/#2 rounded to the closest integer vector.
\newcommand{\Rounded}[2]{\ensuremath{\left\lbrack \frac{#1}{#2} \right\rbrack}}

% Medial axis of #1
\newcommand{\MA}[1]{\ensuremath{\mathrm{MA}(#1)}}
% reach of #1
\newcommand{\Reach}[1]{\ensuremath{\mathrm{reach}(#1)}}
% projection onto #1
\newcommand{\ProjX}[1]{\ensuremath{\pi^{#1}}}




% Digital frontier at step $h$.
%\newcommand{\DFr}[1]{\ensuremath{\Delta_{#1}}}
%\newcommand{\hBd}[1]{\ensuremath{\partial_{#1}}}
\newcommand{\BdZ}[1]{\ensuremath{Bd({#1})}}

%%ParameterFree
% \newcommand{\PCurvHatStar}{\ensuremath{\hat{\Curv}^{*}}}
% \newcommand{\PMeanCurvHatStar}{\ensuremath{\hat{H}^{*}}}


% -----------------------------------------------
% Macro conçu par Jean-David & Olivier Génevaux (permet de citer des papiers dans des images) :
\newcommand\RefGraphicn[3]{\includegraphics[#1]{#2}}

\newcommand\RefGraphics[3]
{%
	\begin{tikzpicture}
		\node
		[
			inner sep=0,
			outer sep=0,
		]
		at (0,0) {\includegraphics[#1]{#2}};

		%\draw [red] (current bounding box.south west) rectangle (current bounding box.north east);

		\node
		[
			xshift=-2mm,
			yshift=1mm,
			above left,
			rectangle,
			draw,
			fill=white!00,
			opacity=.75,
			%font=\sffamily,
		]
		at (current bounding box.south east) {\phantom{#3}};

		\node
		[
			xshift=-2mm,
			yshift=1mm,
			yshift=-0.1mm,
			above left,
			rectangle,
			%font=\sffamily,
		]
		at (current bounding box.south east) {#3};

	\end{tikzpicture}%
}


%%% SVG Images
\newcommand{\svgYes}{~\includegraphics[width=2.6mm]{images/icons/basic14}~}
\newcommand{\svgNope}{~\includegraphics[width=2.6mm]{images/icons/delete85}~}
